% !TEX program = xelatex
% Overleaf: Menu → Compiler → XeLaTeX 로 설정해야 컴파일됨 (pdfLaTeX는 한글에서 타임아웃)
\documentclass[11pt]{article}
\usepackage{kotex}
\usepackage[margin=2.2cm]{geometry}
\usepackage{enumitem}
\usepackage{booktabs}
\usepackage{xcolor}
\usepackage{tcolorbox}
\usepackage{hyperref}
\setlist{itemsep=2pt, topsep=4pt}

\definecolor{codebg}{RGB}{245,245,245}
\newtcolorbox{pathbox}{colback=codebg, colframe=black!25, boxrule=0.4pt, arc=2pt,
  left=6pt, right=6pt, top=4pt, bottom=4pt}

\title{\textbf{프론트엔드 연동 가이드} \\ \large 어디에 뭐가 있는지 --- 7/7 기준}
\author{}
\date{}

\begin{document}
\maketitle
\vspace{-3em}

백엔드 API가 오늘 많이 바뀌어서, 프론트 어느 파일을 고치면 되는지 +
백엔드 어디를 보면 정답이 있는지 위치 기준으로 정리한 문서임.

\section{백엔드에서 봐야 할 위치 (정답이 있는 곳)}

\begin{description}[style=nextline, itemsep=6pt]
  \item[\texttt{docs/specs/백엔드-api-명세.md}]
    모든 엔드포인트의 요청/응답/에러코드 표. 오늘 실제 동작 기준으로 갱신했음.
    헷갈리면 무조건 이 문서가 기준.

  \item[\texttt{backend/server/src/main/java/com/madfinder/server/dto/}]
    응답 JSON의 필드명은 이 폴더의 record 필드명과 1:1로 같음.
    명세랑 코드가 다르게 보이면 dto 파일이 최종 정답.
    \begin{itemize}
      \item \texttt{RecommendResponse.java} --- 추천 응답 (sections.popular/discovery 구조)
      \item \texttt{RecommendStatusResponse.java} --- 정밀모드 폴링 응답 (progress, sections)
      \item \texttt{SimilarGamesResponse.java} --- 비슷한 게임 응답
      \item \texttt{SearchResponse.java} --- 검색 응답
      \item \texttt{GameDetailResponse.java} --- 게임 상세 응답 (screenshots 포함)
      \item \texttt{UserFavoritesResponse.java} --- 닉네임 조회 + 즐겨찾기 응답
      \item \texttt{ErrorResponse.java} --- 에러 공통 형식 \texttt{\{ error, message \}}
    \end{itemize}

  \item[\texttt{backend/server/.../controller/}]
    엔드포인트 경로/메서드가 정의된 곳 (어떤 URL이 있는지 확인용)

  \item[\texttt{backend/config/scoring.json}]
    섹션당 개수(50), alpha 같은 튜닝값. 프론트가 개수를 하드코딩하지 말 것
    (50이 바뀔 수 있으니 배열 길이 그대로 렌더링하면 됨)
\end{description}

\section{프론트에서 고쳐야 할 위치 (작업별)}

\subsection*{작업 1. 추천 응답 sections 구조 반영 \quad (최우선 --- 안 하면 추천 빈 화면)}

\begin{pathbox}
\begin{description}[style=nextline, itemsep=4pt]
  \item[\texttt{frontend/src/types/recommend/index.ts}]
    \texttt{RecommendationsResponse} 타입이 옛 \texttt{\{ recommendations: [...] \}} 형태임.
    \texttt{\{ sections: \{ popular: Item[], discovery: Item[] \} \}} 로 교체.
    Item에 \texttt{genreL2} 필드 추가.
  \item[\texttt{frontend/src/api/recommend/index.ts}]
    응답 파싱이 \texttt{data.recommendations} 를 읽고 있음 $\to$ \texttt{data.sections} 로 교체.
  \item[\texttt{frontend/src/api/recommend/hooks/useRecommendations.ts}]
    위 타입 바뀌면 따라서 수정.
  \item[\texttt{frontend/src/mocks/handlers/recommend.handlers.ts}]
    목업이 \texttt{POST /api/recommendations} (실서버에 없는 경로!) 로 돼있음.
    실서버 경로는 \texttt{POST /api/recommend} 임. 경로랑 응답 형태 둘 다 교체.
  \item[\texttt{frontend/src/mocks/selectors/recommend.selectors.ts}]
    목업 응답 만들어주는 곳. sections 형태로 교체.
  \item[\texttt{frontend/src/pages/recommend/RecommendPage.tsx}]
    두 섹션 UI (``인기 게임'' / ``숨은 발견''). 같은 게임이 양쪽에 있어도 정상.
\end{description}
\end{pathbox}

\subsection*{작업 2. 정밀모드 UI}

새로 만들 것 (기존 파일 없음):
\begin{itemize}
  \item 토글 위치: \texttt{frontend/src/pages/tierlist/TierlistPage.tsx} 의 추천받기 버튼 옆
  \item API 함수: \texttt{frontend/src/api/recommend/} 안에 추가
    \begin{itemize}
      \item \texttt{POST /api/recommend} 바디 \texttt{\{ userId, mode: "precise" \}}
            $\to$ \texttt{\{ jobId \}} 받음
      \item \texttt{GET /api/recommend/status/\{jobId\}} 2~3초 간격 폴링
    \end{itemize}
  \item 진행 화면: status가 running이면 \texttt{progress.current/total/collectingName} 표시.
        done이면 sections 결과 (작업 1과 같은 형태), error면 message 안내
  \item 예외: 409 \texttt{JOB\_ALREADY\_RUNNING} $\to$ ``이미 분석 중'' /
        404 \texttt{JOB\_NOT\_FOUND} $\to$ 재시도 안내
\end{itemize}

\subsection*{작업 3. 상세 페이지 (스크린샷 + 비슷한 게임)}

\begin{pathbox}
\begin{description}[style=nextline, itemsep=4pt]
  \item[\texttt{frontend/src/api/game/index.ts}, \texttt{hooks/useGameDetail.ts}]
    \texttt{GET /api/games/\{id\}} 의 screenshots가 이제 실제 URL 배열로 옴 (캐러셀 재료).
    비슷한 게임용 \texttt{GET /api/games/\{id\}/similar} 호출 함수 추가 ---
    응답 \texttt{\{ similar: [ \{ universeId, name, genreL1, playerCount, iconUrl \} ] \}} 최대 6개.
  \item[\texttt{frontend/src/types/game/index.ts}]
    similar 응답 타입 추가.
  \item[\texttt{frontend/src/pages/game/GameDetailPage.tsx}]
    스크린샷 캐러셀 + 비슷한 게임 섹션. \texttt{videoUrl}은 항상 null이니 무시하고
    영상은 \texttt{GET /api/games/\{id\}/videos} (유튜브 쇼츠)만 쓰면 됨.
  \item[\texttt{frontend/src/mocks/handlers/games.handlers.ts}, \texttt{selectors/}]
    similar 목업 추가 (쓰는 경우).
\end{description}
\end{pathbox}

\subsection*{작업 4. 즐겨찾기 새로고침 버튼}

\begin{pathbox}
\begin{description}[style=nextline, itemsep=4pt]
  \item[\texttt{frontend/src/api/favorites/index.ts}]
    \texttt{GET /api/users/\{username\}/favorites?refresh=true} 로 호출하는 함수/옵션 추가.
  \item[\texttt{frontend/src/pages/favorites/FavoritesPage.tsx} (1페이지)]
    새로고침 버튼: 누르면 refresh=true로 재조회 + 버튼 비활성화 (접속당 1회).
    1~2초 걸릴 수 있으니 로딩 표시, 429(BUSY) 오면 ``잠시 후 재시도'' 안내.
    비활성화 상태는 페이지 새로고침하면 풀려도 됨 (``접속당'' 기준이라).
\end{description}
\end{pathbox}

\subsection*{작업 5. 검색 UI (게임 직접 추가 --- 2페이지)}

\begin{pathbox}
\begin{description}[style=nextline, itemsep=4pt]
  \item[\texttt{frontend/src/api/} 에 검색 함수 추가 (기존 폴더 없음 --- \texttt{api/search/} 새로)]
    \texttt{GET /api/search?q=검색어} ---
    응답 \texttt{\{ results: [ \{ universeId, name, playerCount, iconUrl \} ] \}} 상위 10개.
    로블록스 실시간 호출이라 1초 안팎 걸림. 타이핑마다 부르지 말고
    엔터/검색버튼으로만 호출. 429(BUSY) 오면 ``잠시 후 재시도'' 안내.
  \item[\texttt{frontend/src/pages/tierlist/TierlistPage.tsx}]
    검색 결과에서 게임 선택 $\to$ 티어표에 직접 추가하는 UI.
\end{description}
\end{pathbox}

\section{실행 환경 참고}

\begin{description}[style=nextline, itemsep=6pt]
  \item[\texttt{frontend/.env}]
    \texttt{VITE\_ENABLE\_MSW=false} 면 실서버(localhost:8080) 호출,
    \texttt{true} 면 목업(MSW)으로 개발. 백엔드 안 띄우고 UI만 만들 땐 true로.
  \item[\texttt{frontend/vite.config.ts}]
    \texttt{/api} 요청은 localhost:8080 (로컬 Spring Boot)으로 프록시됨.
\end{description}

실데이터로 보고 싶으면: 7/8 아침에 DB 덤프 보내줄 예정.
그 전까지는 MSW 목업으로 개발하는 게 편함 (단, 목업을 작업 1의 새 형태로
먼저 고쳐야 실서버 붙일 때 화면 코드를 안 건드림).

\vspace{0.5em}
에러 공통 형식 (모든 API 동일):
\texttt{\{ "error": "에러코드", "message": "사람이 읽을 설명" \}} + HTTP 상태코드.
에러코드 전체 표는 \texttt{docs/specs/백엔드-api-명세.md} 하단에 있음.

\end{document}
